\documentclass{article}
\usepackage{graphicx} % Required for inserting images
\usepackage{datetime} % Add this line to include the datetime package
\usepackage{amsmath}
\usepackage{amssymb}
\usepackage{amsthm}
\usepackage{tikz}
%\usepackage{hyperref}
\usepackage{xcolor} % Required for defining custom colors

\definecolor{darkblue}{rgb}{0.0, 0.0, 0.5} % Define a darker blue color

\usepackage[colorlinks=true, linkcolor=darkblue, urlcolor=darkblue, citecolor=darkblue]{hyperref} % Use the darker blue color

% Define theorem style
\theoremstyle{plain} % Bold title, italic body
\newtheorem{theorem}{Theorem}
\newtheorem{lemma}[theorem]{Lemma}
\newtheorem{corollary}[theorem]{Corollary}

\theoremstyle{definition} % Bold title, normal body
\newtheorem{definition}[theorem]{Definition}
\newtheorem{example}[theorem]{Example}

\newcommand{\bZ}{\mathbb{Z}}
\newcommand{\fZ}{\mathbf{Z}}

\title{Homework 8}
\author{Aaron Honculada}
\date{\monthname[\month] \the\year}

\begin{document}

\maketitle
\section*{Problem 1}
\noindent \textit{Show that a group of order 105 is not simple.}

We are given the order of a group $G$ = 105. The prime factorization of $|G| = 3 \cdot 5 \cdot 7$.
When considering the Sylow 3-subgroups, by the Third Sylow Theorem there can only be 1 or 7 Sylow 3-subgroups
since $7 \equiv 1 \equiv 1 \pmod{3}$ and $7 \mid (5 \cdot 7)$. We denote this as $n_3 \in \{1, 7\}$. In the case where
$n_3 = 1$ then by the Second Sylow Theorem, the single Sylow 3-subgroup is normal and therefore $G$ is not simple.
We consider the case where $n_3=7$. If there are seven Sylow 3-subgroups, we know that their intersection is the trivial
group. Therefore there are $7 \cdot (3-1) =14$ elements in $G$ of order 3. 

We now consider the remaining $105 - 15 = 91$ remaining elements. By the Third Sylow Theorem, we know that the number of 
Sylow 5-subgroups are either 1 or 21 since $21 \equiv 1 \equiv 1 \pmod{5}$ and $21 \mid (3 \cdot 7)$. We similarly
denote this as $n_5 \in \{1, 21\}$. By the Second Sylow Theorem, if $n_5 = 1$ then the lone Sylow 5-subgroup
is normal and $G$ is not simple. We consider the case where $n_5 = 21$. The twenty one Sylow 5-subgroups only
intersect at the identy element, therefore we have $21 \cdot (5-1) = 84$ elements of order 5.

We finally consider the remaining $91-84 = 7$ elements of $G$ who do not have order 3 or 5. By the Third
Sylow Theorem, we have it that the number of Sylow 7-subgroups $n_7 \in \{1, 15\}$ since 
$15 \equiv 1 \equiv 1 \pmod{7}$ and $15 \mid (3 \cdot 5)$. By the Second Sylow Theorem, for
$G$ not to be simple it must be that $n_7 = 15$, since if $n_7 = 1$ then the single Sylow 7-Subgroup must
be normal and $G$ is not simple. When we consider the case where $n_7 = 15$ then these
fifteen Sylow 7-subgroups intersect at the trivial group and we have $15 \cdot (7-1) = 90$ elements of order 7.
However this exceeds the remaining number of elements in $G$ since $90 > 7$ and therefore since we cannot have 
$n_3 \neq 1$ and $n_5 \neq 1$ and $n_7 \neq 1$, then one of these three numbers must be equal to 1 and by the Second
Sylow Theorem that sole Sylow p-subgroup must be normal and thus $G$ is not simple.

\section*{Problem 2}
\noindent \textit{Let $N$ be a normal subgroup of a finite group $G$, and let
$K$ be a normal subgroup of $N$.}

\noindent (a) Show that if $K$ is a Sylow p-subgroup of $N$, then $K$ is a normal subgroup of $G$.
\begin{proof}
    We consider $g \in G$. Given that $N \trianglelefteq G$, we know that $gNg^{-1} = N$. Given that $K$ is a Sylow p-subgroup
    of $N$, be definition of Sylow p-subgroup the order of $K$ is the largest prime power $p^n$ that divides $N$. Also given
    that conjugation is a group automorphism that preserves order, we know that $|K| = |gKg^{-1}| = p^n$. Therefore $gKg^{-1} \subseteq
    gNg^{-1}$. Since $gKg^{-1}$ is a subgroup of $N$ with the order $p^n$ we have it that $gKg^{-1}$ is also a Sylow p-subgroup of $N$. By
    the Second Sylow Theorem, if $K$ is normal and a Sylow p-subgroup of $N$, then it must the the only Sylow p-subgroup for that $p$. Therefore
    it must be the case that $gKg^{-1} = K$ for all $g \in G$. This is the definition of normal subgroup and thus when $K$ is a Sylow p-subgroup of $N$, then $K$
    is a normal subgroup of $G$.
\end{proof}

\noindent (b) Give an example showing that the conclusion of part (a) may not hold if $K$ is not a Sylow p-subgroup of $N$.

\section*{Problem 3}


\noindent (a) For any prime $p$, show that if $G$ is a nonabelian group of order $p^3$ with center $\mathbf{Z}$, then $\mathbf{Z} \cong
\mathbb{Z}_p$ and $G/\mathbf{Z} \cong \mathbb{Z}_p \times \mathbb{Z}_p$.
\begin{proof}
    We know from Proposition 3.5 from class lecture that the center of any finite p-group is nontrivial. We have that the order of $\mathbf{Z} \cong
    0 \pmod{p}$ and thus the only values for the order of the center are powers of $p$. Given that $|G| = p^3$ then $|\mathbf{Z}|$ must be either
    $p, p^2, \text{ or, }p^3$. If $|\fZ| = p$ then $\mathbf{Z} \cong \mathbb{Z}_p$ and we are done. 
    
    We now consider the other two cases where $|\fZ| = p^2, p^3$. If $|\fZ| = p^2$ then the quotient 
    $G/\fZ$ has order $p$. This implies that $G/\fZ$ is cyclic which further implies that $G$ is abelian by Proposition 3.6 as presented in lecture. This contradicts the hypothesis that 
    $G$ is nonabelian and therefore $|\fZ| \neq p^2$. If $|\fZ| = p^3$ then the entire group $G$ lies in the center $\fZ$.
    This also implies that $G$ is abelian, which contradicts the hypothesis that $G$ is not abelian. Therefore $|\fZ| = p$ and
    $\fZ \cong \bZ_p$.
    
    When considering the quotient $G/\fZ$, we know that $|G| = p^3$ and $|\fZ|= p$. Therefore we have that
    the order of $G/\fZ = p^2$. By Theorem 3.14, we know that all group of order $p^2$ is abelian and hence is isomorphic to
    $\bZ_{p^2} \cong \bZ_p \times \bZ_p$. Therefore $G/\fZ \cong \bZ_p \times \bZ_p$.
\end{proof}

\noindent (b) Let $H_p$ be a nonabelian group of order $p^3$ consisting of matrices with entries in $\mathbb{Z}_p$ of the form:
\begin{center}
    $\begin{pmatrix}
        1 & a & c \\
        0 & 1 & b \\
        0 & 0 & 1
    \end{pmatrix}$
\end{center}
under multiplication.
Describe the center $\mathbf{Z}$ and give a surjective homomorphism $f: H_p \to \mathbb{Z}_p \times \mathbb{Z}_p$ with kernel $\mathbf{Z}$.

When considering the product of two matrices in $H_p$ we see the following:
\begin{center}
    $\begin{pmatrix}
        1 & a & c \\
        0 & 1 & b \\
        0 & 0 & 1
    \end{pmatrix}$
    $\begin{pmatrix}
        1 & x & z \\
        0 & 1 & y \\
        0 & 0 & 1
    \end{pmatrix}$ =
    $\begin{pmatrix}
        1 & a+x & c+ay+z \\
        0 & 1 & b+y \\
        0 & 0 & 1
    \end{pmatrix}$ 
\end{center}
Therefore multiplication is only commutative when $ay=xb$ for all $x,y \in \bZ_p$. 
This can only be true when both $a=0$ and $b=0$.
We describe $\fZ$, the center of $H_p$, as matrices of the following form
\begin{center}
    $\begin{pmatrix}
        1 & 0 & c \\
        0 & 1 & 0 \\
        0 & 0 & 1
    \end{pmatrix}$
\end{center}
where $c \in \bZ_p$. 

Focusing on the entries in the above matrix multiplication which contain $a+x$ and $b+y$, we 
identify a surjecive homomorphism $f: H_p \to \bZ_p \times \bZ_p$.

Let $f\left(
    \begin{pmatrix}
        1 & a & c \\
        0 & 1 & b \\
        0 & 0 & 1
    \end{pmatrix}\right) = (a,b) \in \bZ_p \times \bZ_p$.

Since 
\begin{center}
    $\begin{pmatrix}
        1 & a & c \\
        0 & 1 & b \\
        0 & 0 & 1
    \end{pmatrix}$
    $\begin{pmatrix}
        1 & x & z \\
        0 & 1 & y \\
        0 & 0 & 1
    \end{pmatrix}$ =
    $\begin{pmatrix}
        1 & a+x & c+ay+z \\
        0 & 1 & b+y \\
        0 & 0 & 1
    \end{pmatrix} \implies f(AB) = (a+x, b+y) = f(A)f(B)$ 
\end{center}
then $f$ is a group homomorphism whose kernel is $\fZ$.

The image of $f$ is all of $\bZ_p \times \bZ_p$ and therefore $f$ is surjective.

\noindent (c) Is $H_2$ isomorphic to the dihedral group $D_4$ or the quaternion group $Q_8$?

We claim that $H_2 \cong D_4$ and not $Q_8$. The elements $\pm i, \pm j, \pm k \in Q_8$ have order 4 and only the
element $-1$ in $Q_8$ has order 2. However in $H_2$ we can consider the two matrices:
\begin{center}
    $\begin{pmatrix}
        1 & 0 & 0 \\
        0 & 1 & 1 \\
        0 & 0 & 1
    \end{pmatrix}$ and 
    $\begin{pmatrix}
        1 & 1 & 0 \\
        0 & 1 & 0 \\
        0 & 0 & 1
    \end{pmatrix}$.
\end{center}
These are two distinct elements in $H_2$ that have order 2 since squaring each matrix
gives the identity matrix. Also each of these matrices
can be mapped to a rotation element of $D_4$. Since $H_2$ has more than one
element that has order two, we claim that $H_2 \not \cong Q_8$ and instead
$H_2 \cong D_4$.

\section*{Problem 4}
\noindent \textit{Let $N$ be a group and Aut($N$) the group of automorphisms of $N$. Let $H$ be a group with a homomorphism $
\varphi: H \to \text{Aut}(N)$ sending $h \to \varphi_h$ for $h \in H$.}

\noindent (a) Show that the set of ordered pairs $(n, h)$ for $n \in N$ and $h \in H$ under the product
\begin{center}
    $(n,h)(n',h') = (n\varphi_h(n'), hh')$
\end{center}
forms a group.

\begin{proof}
We verify that the set of order pairs $(n, h)$ where $n \in N$ and $h \in H$ under
this product satisfy the group axioms: closure, identity, inverse, and associativity.

\begin{itemize}
    \item \textbf{Closure} Given $\varphi_h \in$ Aut$(N)$ and $\varphi_h(n') \in N$ and therefore
    $n \cdot \varphi_h(n') \in N$ and $hh' \in H$ we see that $(n\cdot \varphi_h(n'), hh') \in G$ for all $n,n',h,h' \in G$.
    \item \textbf{Identity} We claim the identity element in $G$ be $(e_N, e_H)$. Given that $\varphi_{e_H}$ is the identity permutation
    and $\varphi_h(e_N) = e_N$, then we see that left multiplying $(e_N, e_H)(n,h) =
    (e_N \cdot \varphi_{e_H}(n), e_Hh)=(n,h)$. Similarly right multiplying we have 
    $(n,h)(e_N, e_H) = (n\cdot \varphi_h(e_N), he_H) = (n, h)$.
    \item \textbf{Inverse} We verify that $(n,h)^{-1} = (\varphi_{h^{-1}}(n^{-1}), h^{-1})$.
        \begin{align*}
            (n,h)(\varphi_{h^{-1}}n^{-1}, h^{-1}) &= (n\cdot \varphi_h(\varphi_h^{-1}(n^{-1})), hh^{-1}) \\
             &= (n\cdot n^{-1},e_H) \\
             &= (e_N, e_H)
        \end{align*}
    \item \textbf{Associativity} Since $\varphi_{h} \in \text{Aut}(N)$ is a group homomorphism, we have: 
        \begin{align*}
    ((n_1,h_1)(n_2,h_2))(n_3,h_3) &= (n_1,h_1)((n_2,h_2)(n_3,h_3)) \\
    (n_1\cdot \varphi_{h_1}(n_2),h_1h_2)(n_3,h_3) &= (n_1,h_1)(n_2 \cdot \varphi_{h_2}(n_3),h_2h_3) \\ 
    (n_1 \cdot \varphi_{h_1}(n_2)\cdot \varphi_{h_1h_2}(n_3), h_1h_2h_3) &= (n_1 \cdot \varphi_{h_1}(n_2)\cdot \varphi_{h_1h_2}(n_3), h_1h_2h_3)
        \end{align*}
\end{itemize}
\end{proof}

\noindent (b) Explain how this can be used to construct a non-abelian group of order $pq$, where $p, q$ are primes with $q \equiv 1 \pmod{p}$.

We can construct a nonabelian group using this product with the cyclic groups $\bZ_p$ and $\bZ_q$. We shall
notate this product as the symbol ``$\rtimes$" and call it semi-direct product. First we identify a bijection between
Aut$(\bZ_q)$ and $\bZ_{q-1}$. For every automorphism $\varphi \in \text{Aut}(\bZ_q)$ we can map it to an integer in $\bZ_{q-1}$. 

\noindent (c) When $N=\bZ_3$ and $H=\bZ_4$, one can use this construction for form
a nonabelian group $T$ of order 12. Describe the group multiplication on ordered pairs $(a,b)$ for
$a \in \bZ_3$ and $b \in \bZ_4$, and show that $T$ is not isomorphic to $A_4$ or $D_6$.

\end{document}